\section{Описание практической части}
\label{sec:Chapter4} \index{Chapter4}

Раздел документирует программную реализацию решения, разработанного
в предыдущей главе, и приводит результаты экспериментов.
Сначала описывается, куда именно в стек DagorEngine встроен метод,
затем --- детали алгоритмического модуля и, наконец, измеренное
ускорение трассировки. При этом все исходные коды, скрипты и сырые материалы
замеров размещены в репозитории работы~\cite{thesisRepo}.

Чтобы было понятно место метода в общей системе, ниже
приведён стек подсистем: от аппаратного
API DirectX~12 до измеряемого прохода \texttt{rtsm::do\_trace}.
В этом стеке предлагаемая предобработка вклинивается между подсистемой
\texttt{bvh} и вызовом построения BLAS.

\begin{figure}[h!]
\centering
\includegraphics[width=0.68\textwidth]{diagrams/scheme_stack.png}
\caption{Стек подсистем: предобработка MERGE-ALL встроена на уровне
семпла \texttt{testGI} перед вызовами \texttt{bvh::add\_mesh} и не
требует изменений в драйвере или ассет-пайплайне. Измеряемая метрика ---
время верхнего прохода \texttt{rtsm::do\_trace}.}
\label{fig:stack}
\end{figure}

\subsection{Архитектура реализации}

Реализация локализована в семпле \texttt{testGI} движка
DagorEngine~\cite{dagorEngine} и состоит из двух частей:

\paragraph{(A) Алгоритмическая часть.} Код алгоритма вынесен в
namespace \texttt{testgi\_dynrend\_bvh\_pre}. Он расположен в файле
\texttt{samples/testGI/prog/test\_app.cpp} (около 600 строк
C++17). Реализация включает чтение треугольников из исходных
\texttt{ShaderMesh::RElem}, глобальную Morton-сортировку по
центроидам, разбиение на пространственно-плотные кластеры,
выделение для каждого кластера собственных vertex/index-буферов и
регистрация их в BVH-подсистеме движка через
\texttt{bvh::add\_mesh()} и \texttt{bvh::add\_instance()}.

\paragraph{(B) Конфигурационная часть.} Опция
\texttt{bvhMergeAllRelems} в файле
\texttt{samples/testGI/game/settings.blk} --- единственный
пользовательский управляющий параметр. Целочисленное значение
задаёт минимальное желаемое число пространственных кластеров на
модель; нулевое или отрицательное значение отключает метод и
возвращает поведение к baseline.

\subsection{Реализация алгоритмического модуля}

\subsubsection{Ключевые структуры данных}

\begin{verbatim}
struct SubMesh { uint32_t triStart, triCount; BSphere3 bsphere; };
struct MergedResourceState {
    uint32_t clusterCount;            // итоговое число кластеров K
    Tab<Sbuffer*> vbs, ibs;           // K vertex- и 16-bit index-буферов
};
struct ClusterRegistry {
    WinCritSec cs;
    eastl::vector_map<uint64_t, MergedResourceState> merged;
};
\end{verbatim}

\texttt{ClusterRegistry} живёт как singleton (ключ --- \texttt{bvh\_id}
ресурса); критическая секция обеспечивает потокобезопасность доступа
из BVH-обновлений.

\subsubsection{Основные функции}

\begin{itemize}
\item \texttt{gather\_relem\_triangles()} читает позиции
  треугольников одного \texttt{RElem}-а (формат \texttt{VSDT\_FLOAT3},
  16-битные индексы; RElem-ы с иным форматом пропускаются с
  предупреждением).
\item \texttt{try\_register\_merged\_resource()} --- основная
  функция: собирает глобальный пул треугольников всех \texttt{RElem}-ов,
  выполняет Morton-сортировку, согласует $K$ с лимитом 16-битного
  индекса, выделяет на каждый кластер свои vertex/index-буферы и
  регистрирует их через \texttt{bvh::add\_mesh()}.
\item \texttt{add\_dynrend\_resource\_to\_bvh()} и
  \texttt{add\_dynrend\_instance\_to\_bvh()} --- модифицированные
  точки входа: при $K\!>\!0$ стандартный per-\texttt{RElem} путь
  заменяется на $K$ кластеров/инстансов с общей трансформацией
  (уникальный 64-битный \texttt{mesh\_id} кодирует пару
  \texttt{(bvh\_id, cluster\_ix)}).
\end{itemize}

\subsubsection{Аппаратные требования к буферам}

Index-буферы кластеров создаются с явным набором флагов:
\begin{verbatim}
const uint32_t ibFlags = SBCF_BIND_INDEX
                       | SBCF_BIND_SHADER_RES
                       | SBCF_MISC_ALLOW_RAW;
Sbuffer *ib = d3d::create_ib(ibSize, ibFlags, "bvh_merged_ib");
\end{verbatim}

Флаги \texttt{SBCF\_BIND\_SHADER\_RES} и
\texttt{SBCF\_MISC\_ALLOW\_RAW} необходимы для того, чтобы
compute-шейдер \texttt{BVH::IndexProcessor} в Dagor мог
использовать буфер индексов как Shader~Resource~View. Без них
runtime срабатывает с \texttt{G\_ASSERT} в начале
BVH-постройки. Это требование выяснено эмпирически по дампам
крашей.

\subsubsection{Согласование $K$ с лимитом 16-битного индекса}

Реализация автоматически поднимает запрошенное $K_{req}$ до
минимально допустимого $K_{min}$. Для целевой модели (около $1{,}22 \cdot 10^6$
треугольников) при запросе \texttt{bvhMergeAllRelems:i = 16}
эффективное K=56 --- именно столько кластеров и регистрируется
(фиксируется в логе инициализации).

\subsection{Конфигурация семпла}

Единственный управляющий параметр --- строка
\texttt{bvhMergeAllRelems:i} в секции \texttt{graphics} файла
\texttt{samples/testGI/game/settings.blk} ($0$ --- метод отключён).
Опция читается при инициализации BVH-контекста и не требует
перезапуска приложения: смена значения вызывает штатный
\texttt{request\_full\_rebuild()}, как при смене RT-preset.

\subsection{Тестовый стенд}

\subsubsection{Сцена}

Сцена представляет собой $4\!\times\!4$ сетку из $N\!=\!16$
инстансов модели \texttt{House, Ko Bulon Don Urak Lawoi Village, Thailand} с шагом 15~м. Между моделями
оставлены проходы для имитации улицы; источник света --- одно
направленное солнце под наклоном $45^\circ$.

\subsubsection{Два варианта разбиения одной модели}

Эксперименты проводились на одной и той же геометрии дома в двух конфигурациях экспорта,
введённых в постановке задачи: \texttt{bad\_house} --- «плохой»
случайный split и
\texttt{bad\_house\_new\_mesh} --- реалистичный материальный split. Для каждого
варианта выполнена серия замеров при разных значениях
\texttt{bvhMergeAllRelems} ($K_{req}$) при фиксированной камере
и одинаковых настройках RTSM.

\subsubsection{Условия замера}

\defitem{Платформа}{Windows~11~x64, DirectX~12, GPU
  GeForce RTX-класса.}
\defitem{Семпл}{\texttt{testGI} из репозитория DagorEngine, сборка
  через штатный \texttt{jam}-скрипт.}
\defitem{Конфигурация}{\texttt{enableBVH=yes}, RTSM включён для
  солнечного источника, остальные RT-эффекты выключены.}
\defitem{Камера}{зафиксирована в одной позиции, обзор сетки
  инстансов под углом сверху-сбоку.}
\defitem{Замер}{медианное GPU-время GPU-маркера
  \texttt{rtsm::do\_trace} за 100 кадров после установления
  тёплого состояния денойзера.}
\defitem{Что не входит в замер}{время сборки BVH
  (\texttt{bvh::build}) и обновления TLAS --- эти операции
  происходят однократно при загрузке сцены и не влияют на
  per-frame стоимость.}

\subsection{Результаты}

\textbf{Главный результат:} предложенный метод снижает время
GPU-трассировки теней \texttt{rtsm::do\_trace} в \textbf{8{,}6 раза}
на синтетическом худшем разбиении (с 9,15 до 1,07~мс)
и в \textbf{2{,}7 раза} на реалистичном разбиении по материалам
(с 2,75 до 1,01~мс) --- при попиксельно идентичной картинке.
Оба результата превышают целевой критерий ($\ge\!2\times$). Ниже
этот результат разбирается подробно: сначала зависимость от числа кластеров
$K_{req}$, затем сводное сравнение и качественная картина
Acceleration~Structure.

\subsubsection{Зависимость времени трассировки от $K_{req}$}

Сводная таблица суммирует sweep по $K_{req}$ для
обоих вариантов разбиения. При $K_{req}$ не менее 64
эффективное число кластеров совпадает с запрошенным (лимит
16-битного индекса для данной модели --- $K_{min}$=56).

\begin{table}[h!]
\centering
\caption{Зависимость \texttt{rtsm::do\_trace} от $K_{req}$ для двух вариантов разбиения модели}
\label{tab:k_sweep}
\small
\begin{tabular}{r|rr|rr}
\hline
& \multicolumn{2}{c|}{\textbf{Случайный split}} &
  \multicolumn{2}{c}{\textbf{Материальный split}} \\
$K_{req}$ &
  \textbf{мс} & \textbf{ускор.} &
  \textbf{мс} & \textbf{ускор.} \\
\hline
0 (baseline) & 9{,}15 & \aptimes{1{,}0} & 2{,}75 & \aptimes{1{,}0} \\
64           & 1{,}07 & \aptimes{8{,}6} & 1{,}01 & \aptimes{2{,}7} \\
96           & 1{,}13 & \aptimes{8{,}1} & 1{,}06 & \aptimes{2{,}6} \\
128          & 1{,}22 & \aptimes{7{,}5} & 1{,}06 & \aptimes{2{,}6} \\
192          & 1{,}23 & \aptimes{7{,}4} & 1{,}11 & \aptimes{2{,}5} \\
256          & 1{,}30 & \aptimes{7{,}0} & 1{,}18 & \aptimes{2{,}3} \\
\hline
\end{tabular}
\end{table}

На «плохом» разбиении метод даёт до \textbf{8{,}6-кратного}
ускорения при $K_{req}$=64: baseline $9{,}15$~мс
против $1{,}07$~мс. На реалистичном разбиении baseline
уже ниже ($2{,}75$~мс), но Morton-кластеризация всё равно снижает
время до $1{,}01$~мс (\aptimes{2{,}7}). В обоих случаях наблюдается
плато $K_{req}$=64…128 и ухудшение при
$K_{req}$=256 из-за роста числа BLAS-инстансов в TLAS
($16$ моделей $\times$ $K$ кластеров). Эта зависимость наглядно
видна на графике ниже.

\begin{figure}[H]
\centering
\includegraphics[width=0.78\textwidth]{diagrams/05_k_sweep_chart.pdf}
\caption{Зависимость времени \texttt{rtsm::do\_trace} от запрошенного
числа кластеров $K_{req}$ для двух вариантов разбиения модели.}
\label{fig:k_sweep}
\end{figure}

\subsubsection{Сводное сравнение baseline и оптимума}

\begin{table}[h!]
\centering
\caption{Baseline и оптимальная конфигурация
($K_{req}$=64) для двух вариантов разбиения}
\label{tab:rtsm_runtime}
\small
\begin{tabular}{lrrr}
\hline
\textbf{Вариант модели} &
  \textbf{baseline, мс} &
  \textbf{$K_{req}$=64, мс} &
  \textbf{Ускорение} \\
\hline
Случайный split (\texttt{bad\_house}) &
  9{,}15 & 1{,}07 & \aptimes{8{,}6} \\
Материальный split (\texttt{bad\_house\_new\_mesh}) &
  2{,}75 & 1{,}01 & \aptimes{2{,}7} \\
\hline
\end{tabular}
\end{table}

Метод даёт наибольший выигрыш там, где исходное разбиение создаёт
сильные перекрытия AABB в TLAS; на реалистичном
разбиении этот эффект меньше (\aptimes{2{,}7}), но остаётся значимым, а оптимум
$K_{req}$ тот же.

\subsubsection{Качественная картина Acceleration Structure}

На рисунках ниже
приведены захваты Acceleration~Structure из Nsight~Graphics
(режим \texttt{Color By: Instance}, каждый цвет --- отдельный
BLAS-инстанс).

\begin{figure}[H]
\centering
\includegraphics[width=0.72\textwidth]{no_optimized.png}
\caption{Случайный split. Наблюдается «шумовая»
картина из перекрывающихся BLAS-инстансов: на любой точке экрана
проецируется множество разных BLAS. Медианное \texttt{rtsm::do\_trace} = \textbf{9{,}15~мс}.}
\label{fig:dagor_with}
\end{figure}

\begin{figure}[H]
\centering
\includegraphics[width=0.72\textwidth]{with_optimization.png}
\caption{После применения метода.
Вместо шума отчётливо видны крупные сплошные цветные пятна --
каждый цвет соответствует одному пространственно-плотному
BLAS-кластеру. В каждом пикселе конкурируют 1--2 BLAS вместо
десятков. \texttt{rtsm::do\_trace} = \textbf{1{,}07~мс}.}
\label{fig:dagor_without}
\end{figure}

Визуальное сравнение подтверждает, что ускорение на «плохом»
разбиении обусловлено именно увеличением пространственной
локальности BLAS-инстансов, а не побочным эффектом подсистемы
RTSM.

Для контраста на отдельном рисунке показан тот же реалистичный
материальный split в исходном состоянии:
модель вручную разбита на группы по материалам, и каждый материал уже 
занимает сравнительно компактную область. Поэтому ещё до применения
метода Acceleration~Structure выглядит локализованной --- мелкие
пространственно-связные цветные пятна, а не сплошной «шум»
перекрывающихся инстансов. Именно это объясняет существенно меньшее
baseline-время на данном варианте ($2{,}75$~мс против $9{,}15$~мс на
случайном split) и более скромный, но устойчивый выигрыш метода
(\aptimes{2{,}7}).

\begin{figure}[H]
\centering
\includegraphics[width=0.72\textwidth]{real_mesh.png}
\caption{Реалистичный материальный split в исходном состоянии. Каждый цвет --- отдельный
материальный BLAS-инстанс. В отличие от случайного split, материалы уже
занимают сравнительно компактные области, поэтому картина выглядит
локализованной ещё \emph{до} применения метода. Медианное
\texttt{rtsm::do\_trace} = \textbf{2{,}75~мс}.}
\label{fig:real_mesh}
\end{figure}

\subsubsection{Выбор $K_{req}$}

Проход по сводной таблице показывает, что при
$K_{req}$=64 достигается минимальное время
\texttt{rtsm::do\_trace} для обоих вариантов модели.
Дальнейший рост~$K$ не улучшает результат: на
$K_{req}$=256 время растёт до $1{,}30$~мс (случайный
split) и $1{,}18$~мс (материальный split) --- накладные расходы
TLAS-обхода начинают доминировать над выигрышем от более мелких
AABB.

\subsection{Обсуждение}

\subsubsection{Интерпретация ускорения}

8{,}6-кратное ускорение на случайном split объясняется через
SAH-стоимость TLAS. В
baseline-конфигурации \texttt{bad\_house} каждая случайная группа
«размазана» по всей площади дома, её \texttt{RElem} образует BLAS с
AABB размером со всю модель, и при 16 инстансах TLAS заполняется
огромными перекрывающимися листьями. Любой луч проверяет пересечение с
большинством этих кандидатов, доля «полезной работы» мала из-за ложных пересечений с AABB --- это синтетический худший
случай, задающий верхнюю границу эффекта.

На реалистичном материальном split baseline уже составляет $2{,}75$~мс:
треугольники частично локализованы по материалам, перекрытия AABB
слабее, и выигрыш меньше (\aptimes{2{,}7}). Именно этот вариант
соответствует реальному пайплайну движка и напрямую подтверждает
гипотезу работы.

После применения метода с $K_{req}$=64 ассет разбивается
на $64$ пространственно-компактных кластера. В TLAS для всей сцены
присутствует $16 \times 64 = 1024$ TLAS-инстансов, ссылающихся на 64 уникальных BLAS, но каждый из инстансов имеет малый AABB,
слабо перекрывающий соседей. Алгоритм обхода TLAS быстро отбрасывает
нерелевантные ветви; при этом при $K_{req}$=256 перекрытия уже
не доминируют, и накладные расходы обхода большого TLAS снова растут.

\subsubsection{Цена ускорения}

Метод увеличивает число BLAS-инстансов на одну модель с
$R \approx 10$ (исходные \texttt{RElem}-ы) до
$K_{req}$ (при оптимуме $K_{req}$=64 --- до
$64$ кластеров). Дополнительный объём памяти на vertex/index-буферы
кластеров составляет $\sim 30$~Мб для всей сцены из 16
инстансов; объём BLAS-узлов растёт пропорционально числу
BLAS-инстансов. Время одноразовой инициализации
(\texttt{bvh::build}) увеличивается на десятки миллисекунд за
счёт большего числа BLAS-построек, но per-frame стоимость
\texttt{rtsm::do\_trace} на «плохом» split падает более чем
в $8$ раз --- что и является целью.

\subsubsection{Ограничения текущей реализации}

Метод имеет ряд ограничений: он теряет per-triangle материальную
атрибуцию, работает только с основным уровнем
детализации, требует ручного выбора $K_{req}$ и неприменим к
анимированным ассетам (скелет, морф-таргеты), вершины которых меняются
каждый кадр. Пути их снятия рассмотрены в заключении.
